% TOP_PROBLEM: {"id":30007636,"problem_number":"LOCAL-30007636","title":"Collective-bargaining redesign","category_id":21,"category":{"id":21,"name":"economics","display_name":"Economics","description":"Open economic theory statements, published research agendas, and proposed scoped specifications; question provenance and historical priority are qualified per record.","slug":"economics","order_index":21,"created_at":"2026-10-08T00:00:00Z"},"status":"research_question","external_url":null,"legacy_ids":[],"published":false,"statement":"In the explicitly specified setting: Which forms of collective bargaining raise worker pay and security while maintaining productivity, investment, employment, and representation of workers outside existing unions? The mathematical statement above fixes the target, assumptions and scope of this catalogue formulation.","economics":{"original_id":125,"field":"Labor markets and work","kind":"design","formalization_basis":"adapted_research_question","source_status":"research_agenda","priority_status":"earliest_found","priority_note":"Earliest directly inspected qualifying passage in this audit; earlier origins were not established. A review or research-agenda statement does not imply original priority.","earliest_verified_reference":{"authors":["Ethan Kaplan","Suresh Naidu"],"title":"Between Government and Market: The Political Economics of Labor Unions","year":2025,"url":"https://econweb.umd.edu/~kaplan/PE_of_Labor.pdf","doi":"10.1146/annurev-economics-051520-025251","locator":"Section 7, Conclusion and Future Research Questions, printed pp. 389–390, bullets on direct policy effects and coverage versus membership","evidence_summary":"The research agenda asks how unions change policy and whether coverage and membership have different political effects; policy optimization is our extension."},"current_reference":{"authors":["Ethan Kaplan","Suresh Naidu"],"title":"Between Government and Market: The Political Economics of Labor Unions","year":2025,"url":"https://econweb.umd.edu/~kaplan/PE_of_Labor.pdf","doi":"10.1146/annurev-economics-051520-025251","locator":"Section 7, Conclusion and Future Research Questions, printed pp. 389–390, bullets on direct policy effects and coverage versus membership","evidence_summary":"The research agenda asks how unions change policy and whether coverage and membership have different political effects; policy optimization is our extension."},"formalization_note":"The cited passage establishes a research direction, not this exact mathematical statement. The notation, population, policy menu, assumptions, and estimands are an original adaptation for this catalogue; no universal unsolved theorem or source priority is claimed."},"statusQualification":"Published question or research agenda verified at the cited locator. The mathematical specification is adapted for this catalogue and includes additional assumptions; source publication does not establish first-ever priority or certify this exact model remains unsolved. The cited passage establishes a research direction, not this exact mathematical statement. The notation, population, policy menu, assumptions, and estimands are an original adaptation for this catalogue; no universal unsolved theorem or source priority is claimed.","statusReviewedAt":"2026-10-08"}
% FIRST_POSTED: 2026-10-08
% ENTRY_KIND: statement_only
% !TeX program = lualatex
\documentclass[11pt]{article}
\usepackage{fontspec}
\usepackage[a4paper,margin=25mm]{geometry}
\usepackage{amsmath,amssymb,mathtools}
\usepackage{xurl}
\usepackage[hidelinks]{hyperref}
\newcommand{\sref}[2]{\hyperlink{source:#1:#2}{[S#2]}}
\setlength{\emergencystretch}{3em}
\begin{document}
% problemId:30007636
\section*{Collective-bargaining redesign}\label{30007636}
\subsection{Definitions and mathematical statement}
Consider a specified economy with workers, firms, and a finite set \(\mathcal B\) of bargaining arrangements. An arrangement \(b\) defines bargaining coverage, representation rules, agreement extension, contract enforcement, and negotiation level. Let \(W_T(b)\), \(L_T(b)\), \(K_T(b)\), and \(R_T(b)\) denote equilibrium compensation, employment, investment, and a preregistered representation index at horizon \(T\). Let \(U_i(b)\) be worker \(i\)'s lifetime utility, measured with an explicit wage--amenity preference model, and let \(\omega_i\ge0\) be declared welfare weights summing to one. Define \(S(b)=E[\sum_i\omega_iU_i(b)]\). The design task is to characterize arrangements maximizing \(S(b)\) subject to specified feasibility, fiscal, and employment constraints. Estimate effects separately for represented workers, uncovered workers, entrants, and consumers. A credible analysis must allow bargaining to change firm entry, productivity, wage-setting, political influence, and who becomes represented. State the variation that distinguishes these mechanisms and the equilibrium closure used in counterfactuals. Existing evidence about a particular union's wage premium cannot identify this entire policy ordering. The optimum is conditional on stated objectives and institutional context, rather than a universally best bargaining system.

\par\textbf{Formulation and scope.} The cited passage establishes a research direction, not this exact mathematical statement. The notation, population, policy menu, assumptions, and estimands are an original adaptation for this catalogue; no universal unsolved theorem or source priority is claimed.

\par\textbf{Open-question provenance.} An explicit question or research agenda was verified at the source locator. This does not by itself certify that every later version remains unresolved \sref{30007636}{1}.

\par\textbf{Historical priority.} Earliest directly inspected qualifying passage in this audit; earlier origins were not established. A review or research-agenda statement does not imply original priority.

\subsection{Short English statement}
In the explicitly specified setting: Which forms of collective bargaining raise worker pay and security while maintaining productivity, investment, employment, and representation of workers outside existing unions? The mathematical statement above fixes the target, assumptions and scope of this catalogue formulation.

\subsection{Sources}
\begin{itemize}\raggedright
\item[\textnormal{[S1]}]\hypertarget{source:30007636:1}{} Ethan Kaplan, Suresh Naidu. 2025. Between Government and Market: The Political Economics of Labor Unions. Section 7, Conclusion and Future Research Questions, printed pp. 389–390, bullets on direct policy effects and coverage versus membership.
\url{https://econweb.umd.edu/~kaplan/PE_of_Labor.pdf}.
DOI: 10.1146/annurev-economics-051520-025251.
{\small\textit{Earliest verified question/agenda reference; historical priority qualified below. The research agenda asks how unions change policy and whether coverage and membership have different political effects; policy optimization is our extension.}}
\end{itemize}
\end{document}
